% 大容量假设空间、插值解集合与优化算法的选解偏差。
\begin{tikzpicture}[figure picture,foundation picture,x=1mm,y=1mm,font=\figurefont\footnotesize]
  \fill[FoundationInput!7] (0,0) ellipse (49 and 32);
  \draw[FoundationInput,line width=1.1pt] (0,0) ellipse (49 and 32);
  \node[font=\figurefont\footnotesize,text=FoundationInput] at (0,27) {架构可表示的假设空间};

  \draw[FoundationLoss,line width=1.1pt]
    plot[smooth cycle,tension=0.8] coordinates {(-31,-7) (-15,10) (5,13) (28,4) (17,-14) (-8,-18)};
  \node[font=\figurefont\footnotesize,text=FoundationLoss] at (-3,-23) {经验风险为零的插值解};

  \draw[FigureMuted,dashed,->] (-45,25) .. controls (-31,19) and (-27,8) .. (-18,4);
  \draw[FigureMuted,dashed,->] (-45,18) .. controls (-25,11) and (-12,5) .. (0,2);
  \draw[FigureMuted,dashed,->] (-45,11) .. controls (-22,4) and (6,0) .. (17,-3);
  \fill[FoundationModel] (0,2) circle (2.2);
  \node[font=\figurefont\footnotesize,text=FoundationModel,align=left,anchor=west] at (5,16)
    {算法输出};
  \draw[foundation flow,FoundationModel] (7,14) -- (1.5,4);

\end{tikzpicture}%
